% Section 8 -- Discussion. GPU and precision paragraph from Efe (was 08_discussion_efe.tex).
% Numbers come from Section 7; nothing new is measured here.

\section{Discussion}
\label{sec:discussion}

\paragraph{Strengths and weaknesses of QAOA.}
Good angles are similar across instances and sizes, so transferred angles beat INTERP at lower cost (E2) and carry QAOA to $n = 28$ with no optimisation (E5), which matters most on real hardware, where every evaluation is slow and noisy.
QAOA's main cost is the optimisation loop: hundreds of evaluations, against milliseconds for CP-SAT (E3).
$P(\text{opt})$ also falls with $n$, so larger problems need more shots, and noise caps the useful depth at $p = 1$--$2$ at a 1\,\% error rate (E6).

\paragraph{GPU and precision.}
Our $2{,}500$-fold GPU speed-up at $n = 20$ is worth only about 11 extra qubits, because each qubit doubles the cost, and one 48\,GB A40 holds at most 32 qubits in fp32.
Single precision was the right choice on the A40.
From 24 qubits on it ran about six times faster than fp64, and in our check at $n = 12$ it never changed the best cost found.

\paragraph{Limitations.}
We schedule onto only two nodes; $k$ nodes would need several qubits per job and constraints to keep each job on exactly one node.
Above $n = 24$ we know only one optimum and its mirror, so $P(\text{opt})$ there is a lower bound.
Our noise model is depolarising only: it leaves out readout errors, decoherence and the SWAP gates that limited qubit connectivity would add, and we do not re-optimise angles under noise.
Finally, our timings are for simulation, not quantum hardware.
